The use of computers is spreading to every field of economic
activity. Most users employ application
soft\-ware
to carry out one or more specific tasks: text
editing, ac\-count\-ing, games, etc. Out of a desire to make the most of the
power of the machine, for the fun of it, or simply out of a
wish to understand how computers work,
many budding computer scientists want to learn
programming, sometimes on their own. The task often turns out to be long, difficult,
even off-putting. We would like to make it a creative activity,
one that trains the mind and is fascinating.

\medskip

Many introductions to programming rely on prior knowledge
of the architecture of computers or use mathematical
notions that are beyond the reach of children or of adults
without a scientific background \cite{LT}\@%
\footnote{
Scientists are used to `citing their sources' by
giving references to the articles in scientific journals and to the
books of the authors whose ideas they use. The modern way to
cite sources is to use a number placed between square brackets
(as for example \cite{LT}) which refers to a list of
references, placed at the end of the book, called the
{\em bibliography} (see page
\pageref{ch:biblio}).}. {\footnotesize (The small
raised numbers refer to a
 footnote explaining the preceding word. For example, the note at
 the bottom of this page explains the use of numbers between square brackets
 such as \cite{LT}.)} We propose to proceed differently.

\medskip

We use drawing, which is within everyone's reach. When the book was first published
in 1991, this approach was
relatively new in the teaching of computer science, since the
strik\-ing progress of technology that brought microcomputers with very
powerful graphic screens within reach of the
general public dated only from the 1980s. It is amusing to note that
the origins of our graphic activity go back to prehistory (more precisely to
the period of cave art%
%
\footnote{Painting on the walls of a cave.
} between 22000 and 13000 BC). Computer graphics, which was born for the general public
in the 1980s, will therefore allow us, 20000 years later, to draw with
the same pleasure as our distant ancestors.

\medskip

This book is for everyone. Nursery-school children are able
to make stylized drawings by hand and, with a little imagination, to
reproduce them on the computer
screen using an interactive drawing program (chapters 1 and 2).
Primary-school pupils can write programs to produce
figurative drawings (chapters 3 to 5). Very quickly they will discover how
eas\-i\-ly they can program the computer to make simple geometric drawings based on
repetitive patterns that they would have neither the patience nor the precision to carry out
by hand (chapters 6 and 7). Middle-school pupils will put their
rudimentary knowledge of algebra and analytic geometry to work to write
drawing programs
using elementary calculations and Cartesian coordinate systems (chapters 8 to 17).
They will find an intuitive
basis for putting into practice
concepts that may seem abstract. High-school
pupils are able to take in more computer-oriented notions such as
coding, or even less elementary notions such as
recursion, which become obvious through drawing (chapters 17 to 22). University students
will enjoy understanding simply the programming that sometimes seems obscure when
it is presented too quickly in the first years of university.

\medskip

Since programming languages do not offer the desirable graphic
possibilities, we have added a li\-brary to the \python\ language\@%
%
\footnote{\ \python\ is a programming language designed by Guido van Rossum
in the Netherlands in the early 1990s. It is used both for teaching programming and for
writing large software. Our programs use \python\ (version~3) and its
graphics library {\tt tkinter}; we have tried them on Macintosh, where the window
of the \robot\ is an ordinary macOS window.}
%
which allows us
to obtain in a simple way effects that are simple (such as arcs of circles) or
spectacular (such as painting) without being tied to the technology of
current graphic screens and without resorting to advanced notions of
programming such as data structures. This library is in
the directory \texttt{LeRobot\_Python\_EN}, and
all the examples and the solution programs of the exercises proposed in the book
are provided in the folder {\tt LeRobot\_Python\_EN/book/}%
%
\footnote{\ The directory \texttt{LeRobot\_Manual\_Python} contains the installation manual
\texttt{LeRobot\_Install.pdf} and a user manual \texttt{LeRobot\_Manual.pdf} of the \robot. The programs of the book are in the folder
{\tt book}, sorted by chapter. To run a program, for example
{\tt Square.py}, just type the command {\tt ./robot.sh book/03\_programming/Square.py}
in a terminal window opened in the folder {\tt LeRobot\_Python\_EN}. Another solution
is to drag (with the mouse, (left) button pressed (do not release it), or with two fingers on the touchpad (trackpad)) the icon of the file containing the program onto the icon of the robot \texttt{LeRobotPythonEN.app}.}%
. It offers many possibilities that are not explored in the book,
such as, for specialists, creating interactive learning programs
like those that are provided for addition, Cartesian coordinates and
translation.

\medskip

This book should be read in successive stages, taking the time and the pleasure
to do many personal experiments on the computer, solving exercises,
modifying the programs proposed as solutions and
looking for new drawing themes.


\vspace{1.0cm}

%%%%%%%%%%%%%%%%%%%%%%%%%%%%%%%%%%%%%%%%%%%%%%%%%%%%%%%%%%%%%%%%%%%%%%%%%%%
\noindent \framebox[\textwidth]{
\begin{minipage}{\myparboxwidth}
\setlength{\parindent}{2em}
\small \sf
\vspace{\abovedisplayskip}
\section*{Addenda}

At the end of each chapter we give, as addenda, additional notes
intended for readers already familiar with computer science.

In this book we introduce the basic activities of programming
such as typing in the text, interpretation, execution
and debugging of programs. The programming notions are presented
in \python: structuring programs into procedures (functions), bounded
iterations, integer expressions, constant declarations,
parameters (passed by value), tests, boolean expressions,
unbounded iterations, coding of information, rational
(floating-point) expressions, recursion, variables and assignments.
Our approach based on drawing gives a graphic intuition
of these notions and avoids the more complex problems of
specification and correctness of programs (which are visually
obvious). Data structures (lists,
tuples, dictionaries and objects in \python) are deliberately left for
later, more advanced learning, but are illustrated graphically.

Learning programming at a very elementary level such as
the one proposed in this book develops a sense of the passage from the concrete
to the abstract, an intuition of formal structures, inductive reasoning and the ability to
formalize, which are at the
basis of programming
but also the starting point of any scientific approach.

\vspace{\abovedisplayskip}
\end{minipage}
}
%%%%%%%%%%%%%%%%%%%%%%%%%%%%%%%%%%%%%%%%%%%%%%%%%%%%%%%%%%%%%%%%%%%%%%%%%%%%
